% ===== Block B11: 原文 p.91–101 | §5.1 射影平坦联络、§5.2 扭曲微分算子层 | 译手 tr-ch4 =====

\subsection{射影平坦联络}

% [原文 p.91]

设 $\mathfrak F=(\pi:\cC\to\cB;s_1,\dots,s_N;\eta_1,\dots,\eta_N)$ 是带形式邻域的 $N$-点曲线族，使得 $\mathfrak F^{(0)}:=(\pi:\cC\to\cB;s_1,\dots,s_N)$ 是 $N$-点曲线的通用族。我们将自由地使用第 4 章的记号，特别是 4.2 中的记号。

由命题 4.7 与命题 4.12 所证，层 $\cL(\mathfrak F)$ 从左作用在 $\cV_{\tilde\lambda}(\mathfrak F)$ 与 $\cV^\dagger_{\tilde\lambda}(\mathfrak F)$ 上。此外，对 $\vec\ell\in\cL(\mathfrak F)$，若 $\theta(\vec\ell)\ne0$，则该作用定义了一个一阶微分算子。本节我们先更仔细地研究 $\theta(\vec\ell)=0$ 时 $\vec\ell$ 的作用。注意由 (4.12) 我们有正合列
\begin{equation}
 0\to\pi_*(\Theta_{\cC/\cB}(*S))\xrightarrow{\ \iota\ }\cL(\mathfrak F)
 \xrightarrow{\ \theta\ }\Theta_\cB(-\log D)\to0. \tag{5.1}
\end{equation}

\begin{lemma}[5.1]
 假设 $\cB$ 足够小，使得我们能找到满足 $\operatorname{Res}^2(\omega)=1$ 的对称双微分 $\omega$。则存在唯一的 $\cO_\cB$-模同态
 \[
  a:\pi_*\Theta_{\cC/\cB}(*S)\to\cO_\cB,
 \]
 它与 $\omega$ 的选取无关，使得对任何 $\vec\ell\in\cL(\mathfrak F)$ ——只要它是 $\pi_*\Theta_{\cC/\cB}(*S)$ 中某个元素经正合列 (5.1) 的像——我们有
 \[
  D(\vec\ell)=a(\vec\ell)\cdot\id
 \]
 作为从左作用于 $\cV_{\tilde\lambda}(\mathfrak F)$ 的线性算子，以及
 \[
  D(\vec\ell)=-a(\vec\ell)\cdot\id
 \]
 作为从左作用于 $\cV^\dagger_{\tilde\lambda}(\mathfrak F)$ 的算子。
\end{lemma}

\noindent\textit{证明.}\quad 由于 $\theta(\vec\ell)=0$，由 (4.17)，对任意 $\langle\Psi|\in\cV^\dagger_{\tilde\lambda}(\mathfrak F)$ 与 $|\Phi\rangle\in\cV_{\tilde\lambda}(\mathfrak F)$ 我们有
\begin{align*}
 \langle\Psi|\{D(\vec\ell)|\Phi\rangle\}
 &=-\sum_{j=1}^{N}\langle\Psi|\rho_j(T[\ell_j])|\Phi\rangle\\
 &=-\sum_{j=1}^{N}\operatorname{Res}_{\xi_j=0}\big(\ell_j(\xi_j)\langle\Psi|T(\xi_j)|\Phi\rangle\,d\xi_j\big)
\end{align*}
其中 $\underline{\ell}_j=\ell_j(\xi_j)\dfrac{d}{d\xi_j}$，并且
\begin{align*}
 &\langle\Psi|T(\xi_j)|\Phi\rangle(d\xi_j)^2\\
 &\quad=\lim_{\zeta_j\to\xi_j}\Big\{\frac{1}{2(\ell+g^*)}
   \sum_{a=1}^{\dim\mathfrak g}\langle\Psi|J^a(\zeta_j)J^a(\xi_j)|\Phi\rangle\,d\zeta_jd\xi_j
   -\frac{c_v}{2(\zeta_j-\xi_j)^2}\langle\Psi|\Phi\rangle\,d\zeta_jd\xi_j\Big\}.
\end{align*}

% [Tutor Note: 这里 $\ell$ 是水平（level），$g^*$ 是对偶 Coxeter 数，故 $\ell+g^*$
%  就是 §4 中反复出现的 $\kappa$（见记号表 $\kappa=g^*+\ell$）；$c_v$ 是真空模
%  （Virasoro 代数作用）的中心荷。第一个式子只是把 $T(\xi_j)$ 的定义按 OPE 写开：
%  第一项是 $\mathfrak g$ 流 $J^a$ 的二次式，第二项是中心荷对应的 $(\zeta_j-\xi_j)^{-2}$
%  极点。下面把 $T$ 换成去掉这一极点的 $\tilde T$，是本节全部计算的关键。]

% [原文 p.92]

选取对称双微分
\[
 \omega\in H^0(\cC\times_\cB\cC,\omega_{\cC\times_\cB\cC/\cB}(2\Delta)),
\]
使其满足 $\operatorname{Res}^2(\omega)=1$。这样的双微分的存在性已在推论 1.13 中证明。令
\begin{align}
 &\langle\Psi|\tilde T(z)|\Phi\rangle(dz)^2 \nonumber\\
 &\quad=\lim_{w\to z}\Big\{\frac{1}{2(\ell+g^*)}
   \sum_{a=1}^{\dim\mathfrak g}\langle\Psi|J^a(w)J^a(z)|\Phi\rangle\,dwdz
   -\frac{c_v}{2}\omega(w,z)\langle\Psi|\Phi\rangle\,dwdz\Big\}. \tag{5.2}
\end{align}

射影联络 $S_\omega(z)(dz)^2$ 由
\[
 S_\omega(z)(dz)^2=6\lim_{w\to z}\Big\{\omega(w,z)dwdz-\frac{dwdz}{(w-z)^2}\Big\}
\]
定义。它不仅依赖于 $\omega$ 的选取，还依赖于局部坐标的选取（见定理 1.15）。由于我们有
\[
 \langle\Psi|T(\xi_j)|\Phi\rangle(d\xi_j)^2
 =\langle\Psi|\tilde T(\xi_j)|\Phi\rangle(d\xi_j)^2
 +\frac{c_v}{12}\langle\Psi|\Phi\rangle S_\omega(\xi_j)(d\xi_j)^2,
\]
我们有
\begin{align}
 \langle\Psi|\{D(\vec\ell)|\Phi\rangle\}
 &=-\sum_{j=1}^{N}\operatorname{Res}_{\xi_j=0}
   \big(\ell_j(\xi_j)\langle\Psi|\tilde T(\xi_j)|\Phi\rangle\,d\xi_j\big)
   \nonumber\\
 &\quad-\frac{c_v}{12}\langle\Psi|\Phi\rangle
   \sum_{j=1}^{N}\operatorname{Res}_{\xi_j=0}
   \big(\ell_j(\xi_j)S_\omega(\xi_j)d\xi_j\big). \tag{5.3}
\end{align}

由于 $\ell_j(z)\langle\Psi|\tilde T(z)|\Phi\rangle dz$ 是整体亚纯 1-形式，且其极点只出现在 $s_j(\cB)$ 处，故上述等式右端的第一项为零。因此，若令
\begin{equation}
 a_\omega(\vec\ell)=-\frac{c_v}{12}\sum_{j=1}^{N}\operatorname{Res}_{\xi_j=0}
 \big(\ell_j(\xi_j)S_\omega(\xi_j)d\xi_j\big), \tag{5.4}
\end{equation}
则我们有
\[
 \langle\Psi|\{D(\vec\ell)|\Phi\rangle\}=a_\omega(\vec\ell)\langle\Psi|\Phi\rangle.
\]
现在 $\cV_{\tilde\lambda}(\mathfrak F)$ 与 $\cV^\dagger_{\tilde\lambda}(\mathfrak F)$ 是局部自由且彼此对偶的，故我们得到
\[
 D(\vec\ell)|\Phi\rangle=a_\omega(\vec\ell)|\Phi\rangle.
\]
下面我们证明 $a_\omega(\vec\ell)$ 与 $\omega$ 的选取无关。若
\[
 \omega'\in H^0(\cC\times_\cB\cC,\omega_{\cC\times_\cB\cC/\cB}(2\Delta))
\]
是另一个满足 $\operatorname{Res}^2(\omega)=1$ 的对称双微分，则 $\omega-\omega'$ 是 $\cC\times_\cB\cC$ 上的全纯双微分。因此，$S_\omega-S_{\omega'}$ 也是 $\omega_{\cC/\cB}^{\otimes2}$ 在 $\cC$ 上的全纯截面。

% [Tutor Note: $\operatorname{Res}^2(\omega)=1$ 这个归一化条件的唯一用处就在这里：两个
%  满足它的双微分之差的双留数为零，于是 $\omega-\omega'$ 在全对角线上没有极点、是
%  全纯的；再加上 $S_\omega$ 的定义中那个 $6\lim$，才有 $S_\omega-S_{\omega'}$ 全纯。
%  否则 $a_\omega(\vec\ell)$ 会差一个非零的留数，从而依赖 $\omega$。]

% [Tutor Note: (5.3) 的推导要点：把 $T=\tilde T+\frac{c_v}{12}S_\omega$ 代入第一式，前一项
%  是 $\tilde T$ 的留数和，后一项提出与态无关的因子 $\frac{c_v}{12}\langle\Psi|\Phi\rangle$。
%  这正是为什么最后 $D(\vec\ell)$ 在 $\cV_{\tilde\lambda}$ 上等于乘以常数 $a_\omega(\vec\ell)$：
%  因为 $\tilde T$ 项的留数和为零（$\ell_j(z)\langle\Psi|\tilde T(z)|\Phi\rangle dz$ 是整体
%  亚纯形式，极点只在标记点 $s_j(\cB)$，而 $\ell_j$ 在 $s_k(\cB)$（$k\ne j$）处没有极点）。
%  另外，按 (5.4) 定义 $a_\omega$ 时要注意 $S_\omega$ 本身不是张量：它在坐标变换下差一个
%  Schwarz 导数项（见定理 1.15），所以 $a_\omega(\vec\ell)$ 一般不是整体定义的函数，
%  只有当 $\vec\ell$ 是 $\iota$ 的像（此时可用 (5.1) 的 $\pi_*\Theta_{\cC/\cB}(*S)$ 部分把它
%  约化掉）时才与坐标无关。这正是引理 5.1 要求 $\vec\ell$ 属于该像的原因。]

% [原文 p.93]

设 $\tau$ 是 $\pi_*\Theta_{\cC/\cB}(*S)$ 中满足 $\iota(\tau)=\vec\ell$ 的元素。则 $\tau(z)(S_\omega(z)-S_{\omega'}(z))dz$ 是 $\cC$ 上的亚纯 1-形式。因此我们有
\[
 \sum_{j=1}^{N}\operatorname{Res}_{\xi_j=0}
 \big(\ell(\xi_j)(S_\omega(\xi_j)-S_{\omega'}(\xi_j))d\xi_j\big)=0.
\]
于是我们得到
\[
 a_\omega(\vec\ell)=a_{\omega'}(\vec\ell).
\]
这就证明了引理的第一部分。由 (4.21) 同理可证第二部分。\hfill$\square$

上述证明还给出下述推论。

\begin{corollary}[5.2]
 在与引理 5.1 相同的假设下，对每个 $\vec\ell\in\cL(\mathfrak F)$，若用 (5.4) 定义 $a_\omega(\vec\ell)$，则
 \[
  a_\omega:\cL(\mathfrak F)\to\cO_\cB
 \]
 是 $\cO_\cB$-模同态。
\end{corollary}

\begin{remark}[5.3]
 在上述推论中，由于 $\vec\ell$ 不必是整体亚纯向量场的像，$a_\omega(\vec\ell)$ 确实依赖于 $\omega$ 的选取。
\end{remark}

现在，只要 $\cB$ 足够小，我们就可以定义 $\cV_{\tilde\lambda}(\mathfrak F)$ 与 $\cV^\dagger_{\tilde\lambda}(\mathfrak F)$ 上的联络了。固定对称双微分 $\omega\in H^0(\cC\times_\cB\cC,\omega_{\cC\times_\cB\cC/\cB}(2\Delta))$，使其满足 $\operatorname{Res}^2(\omega)=1$。对每个元素 $X\in\Theta_\cB(-\log D)$，存在元素 $\vec\ell\in\cL(\mathfrak F)$ 使 $\theta(\vec\ell)=X$。定义从左作用于 $\cV^\dagger_{\tilde\lambda}(\mathfrak F)$ 的算子 $\nabla^{(\omega)}_X$ 为
\begin{equation}
 \nabla^{(\omega)}_X(\langle\Psi|)=D(\vec\ell)(\langle\Psi|)+a_\omega(\vec\ell)(\langle\Psi|), \tag{5.5}
\end{equation}
其中 $\langle\Psi|$ 表示 $\cV_{\tilde\lambda}(\mathfrak F)$ 中的一个元素。

\begin{proposition}[5.4]
 $\nabla^{(\omega)}_X$ 是定义良好的，并且具有下述性质：
 \begin{equation}
  \nabla^{(\omega)}_X(f\langle\Psi|)=X(f)\langle\Psi|+f\nabla^{(\omega)}_X(\langle\Psi|), \tag{5.6}
 \end{equation}
 其中 $f\in\cO_\cB$。因此，对应 $X\mapsto\nabla^{(\omega)}_X$ 定义了 $\cV^\dagger_{\tilde\lambda}(\mathfrak F)$ 上一个沿 $D$ 具有正则奇点的联络。
\end{proposition}

% [Tutor Note: (5.5) 之所以「定义良好」，正是引理 5.1：满足 $\theta(\vec\ell)=X$ 的
%  $\vec\ell$ 只差一个 $\vec\ell-\vec\ell'$（$\theta=0$，即来自 $\pi_*\Theta_{\cC/\cB}(*S)$ 的
%  像），而引理 5.1 说这部分作用恰是乘以 $a(\vec\ell-\vec\ell')$；右边 $+a_\omega(\vec\ell)$
%  正是把它抵消掉。所以「射影平坦」不是附加条件，而是这一构造的必然产物：换 $\vec\ell$
%  代表只改变一个与态无关的 1-形式。]

\noindent\textit{证明.}\quad 取另一个满足 $\theta(\vec\ell')=X$ 的 $\vec\ell'\in\cL(\mathfrak F)$。则我们有
\[
 a(\vec\ell-\vec\ell')=a_\omega(\vec\ell)-a_\omega(\vec\ell').
\]
由引理 5.1 我们有
\[
 (D(\vec\ell)-D(\vec\ell'))\langle\Psi|=\sum_{j=1}^{N}a_\omega(\vec\ell-\vec\ell')\langle\Psi|.
\]
于是，作为 $\cV^\dagger_{\tilde\lambda}(\mathfrak F)$ 上的算子，我们有等式
\[
 D(\vec\ell)+a_\omega(\vec\ell)\cdot\id=D(\vec\ell')+a_\omega(\vec\ell')\cdot\id.
\]
因此 $\nabla^{(\omega)}_X$ 不依赖于 $\vec\ell$ 的选取。由命题 4.6 第 3) 条，(5.6) 成立。\hfill$\square$

\begin{theorem}[5.5]
 $\cV^\dagger_{\tilde\lambda}(\mathfrak F)|_{\cB\setminus D}$ 上的联络 $\nabla^{(\omega)}$ 是射影平坦的。
\end{theorem}

% [Tutor Note: 证明中右端的 $\sum_{j=1}^{N}$ 原书如此，但被加项
%  $a_\omega(\vec\ell-\vec\ell')\langle\Psi|$ 并不依赖 $j$，照字面等于
%  $Na_\omega(\vec\ell-\vec\ell')\langle\Psi|$，与上一式
%  $a(\vec\ell-\vec\ell')=a_\omega(\vec\ell)-a_\omega(\vec\ell')$ 及引理 5.1 的结论不符。
%  正确的读法是去掉求和号（即 $D(\vec\ell-\vec\ell')=a(\vec\ell-\vec\ell')\id$）。]

% [原文 p.94]

\noindent\textit{证明.}\quad 令 $\cB'=\cB\setminus D$。对元素 $X,Y\in\Theta_{\cB'}$，选取 $\vec\ell,\vec m\in\cL(\mathfrak F)$ 使 $\theta(\vec\ell)=X$、$\theta(\vec m)=Y$。为记号简便，证明中我们略去联络记号的上标 $(\omega)$。曲率由
\[
 R(X,Y)=[\nabla_X,\nabla_Y]-\nabla_{[X,Y]}
\]
给出。直接计算可得
\begin{align*}
 \nabla_X(\nabla_Y(\langle\Psi|))
 &=D(\vec\ell)\big(D(\vec m)\langle\Psi|+a_\omega(\vec m)\langle\Psi|\big)
   +a_\omega(\vec\ell)\big(D(\vec m)\langle\Psi|-a_\omega(\vec m)\langle\Psi|\big)\\
 &=D(\vec\ell)D(\vec m)\langle\Psi|+X(a_\omega(\vec m))\langle\Psi|
   +a_\omega(\vec m)(D(\vec\ell)\langle\Psi|)\\
 &\quad+a_\omega(\vec\ell)(D(\vec m)\langle\Psi|)+a_\omega(\vec\ell)a_\omega(\vec m)\langle\Psi|.
\end{align*}
因此我们有
\[
 [\nabla_X,\nabla_Y]\langle\Psi|=[D(\vec\ell),D(\vec m)]\langle\Psi|
 +X(a_\omega(\vec m))\cdot\langle\Psi|-Y(a_\omega(\vec\ell))\cdot\langle\Psi|.
\]
另一方面，由 (4.16)，对 $\theta(\vec n)=[X,Y]$ 我们有
\[
 \vec n=[\vec\ell,\vec m]_d=[\vec\ell,\vec m]_0+X(\vec m)-Y(\vec\ell),
\]
并且
\[
 D(\vec n)=[D(\vec\ell),D(\vec m)]
 -\frac{c_v}{12}\sum_{j=1}^{N}\operatorname{Res}_{\xi_j=0}
 \Big(\frac{d^3\ell_j}{d\xi_j^3}m_jd\xi_j\Big)\cdot\id.
\]
于是我们有
\[
 \nabla_{[X,Y]}\langle\Psi|=[D(\vec\ell),D(\vec m)]\langle\Psi|
 +a_\omega(\vec n)\cdot\langle\Psi|
 -\frac{c_v}{12}\sum_{j=1}^{N}\operatorname{Res}_{\xi_j=0}
 \Big(\frac{d^3\ell_j}{d\xi_j^3}m_jd\xi_j\Big)\cdot\langle\Psi|.
\]
令
\[
 A(\vec\ell,\vec m)=\Big\{X(a_\omega(\vec m))-Y(a_\omega(\vec\ell))-a_\omega(\vec n)
 +\frac{c_v}{12}\sum_{j=1}^{N}\operatorname{Res}_{\xi_j=0}
 \Big(\frac{d^3\ell_j}{d\xi_j^3}m_jd\xi_j\Big)\Big\}\id.
\]
按定义我们有
\[
 a_\omega(\vec m)=-\frac{c_v}{12}\sum_{j=1}^{N}\operatorname{Res}_{\xi_j=0}
 \big(m_j(\xi_j)S_\omega(\xi_j)d\xi_j\big).
\]
于是我们有
\begin{align}
 A(\vec\ell,\vec m)&=\frac{c_v}{12}\Big\{\sum_{j=1}^{N}\operatorname{Res}_{\xi_j=0}
 \big(m_j(\xi_j)\ell_j'(\xi_j)S_\omega(\xi_j)+m_j(\xi_j)X(S_\omega(\xi_j))\,d\xi_j\big)
 \nonumber\\
 &\quad-\sum_{j=1}^{N}\operatorname{Res}_{\xi_j=0}
 \big(\ell_j(\xi_j)m_j'(\xi_j)S_\omega(\xi_j)+\ell_j(\xi_j)Y(S_\omega(\xi_j))\,d\xi_j\big)
 \nonumber\\
 &\quad+\sum_{j=1}^{N}\operatorname{Res}_{\xi_j=0}
 \Big(\frac{d^3\ell_j}{d\xi_j^3}m_jd\xi_j\Big)\Big\}\id \tag{5.7}
\end{align}
其中 $\ell_j'(\xi_j)$ 与 $m_j'(\xi_j)$ 是关于变量 $\xi_j$ 的导数。由于 $A(\vec\ell,\vec m)$ 等于 $R(X,Y)$，故 $A(\vec\ell,\vec m)$ 与 $\vec\ell$、$\vec m$ 到 $\cL(\mathfrak F)$ 的提升的选取无关。因此曲率形式 $R$ 可视为全纯 2-形式，从而该联络是射影平坦的。\hfill$\square$

% [Tutor Note: 直接计算的第一行里，第二项 $a_\omega(\vec\ell)(\cdots)$ 的括号中原文写作
%  $D(\vec m)\langle\Psi|-a_\omega(\vec m)\langle\Psi|$（减号）；但按 (5.5) 应是
%  $D(\vec m)\langle\Psi|+a_\omega(\vec m)\langle\Psi|$，下一行也正是按加号展开的
%  （出现 $+a_\omega(\vec\ell)a_\omega(\vec m)\langle\Psi|$）。照原文排，减号应为笔误。
%  另外，「$R$ 可视为全纯 2-形式 $\Rightarrow$ 射影平坦」的中间一步是：
%  $A(\vec\ell,\vec m)$ 是 $\cO_{\cB'}$ 中的函数乘以 $\id$，故 $R(X,Y)$ 恒是标量算子；
%  而射影平坦的定义正是曲率取值于标量（即 $R(X,Y)=A(X,Y)\id$，$A$ 为 2-形式）。
%  证明 $A$ 与提升无关用的是：换 $\vec\ell\mapsto\vec\ell+\vec\ell_0$（$\theta(\vec\ell_0)=0$）
%  时 $D$ 只差引理 5.1 的标量项，代入 $A$ 的定义恰抵消。原文把这一句压缩成了
%  "Since $A(\vec\ell,\vec m)$ is equal to $R(X,Y)$"。]

% [原文 p.95]

\begin{theorem}[5.6]
 对族 $\mathfrak F$，若我们在每条纤维上选取一个全纯依赖于底参数 $\cB$ 的素形式（prime form），则共形块层 $\cV^\dagger_{\tilde\lambda}(\mathfrak F)$ 带有射影平坦联络，其曲率形式具有形状
 \[
  A(\vec\ell,\vec m)=\frac{c_v}{12}\Big\{\sum_{j=1}^{N}\operatorname{Res}_{\xi_j=0}
  \big((m_j(\xi_j)X(S_\omega(\xi_j))-\ell_j(\xi_j)Y(S_\omega(\xi_j)))d\xi_j\big)\Big\},
 \]
 其中对称双微分 $\omega$ 由素形式定义（见命题 1.12）。
\end{theorem}

\begin{corollary}[5.7]
 若族 $\mathfrak F=(\pi:\cC\to\cB;s_1,\dots,s_N;\eta_1,\dots,\eta_N)$ 是带形式坐标的 $N$-点 $\PP^1$ 族或 $N$-点椭圆曲线族，则共形块层带有平坦联络。
\end{corollary}

\noindent\textit{证明.}\quad 对 $\PP^1$，我们可以选取满足 $\operatorname{Res}^2(\omega)=1$ 的对称双微分
\[
 \omega(x,y)=\frac{dxdy}{(x-y)^2},
\]
其中 $x,y$ 是 $\C\subset\PP^1$ 的坐标。此时射影联络为零，故 $\nabla^{(\omega)}$ 是平坦的。对椭圆曲线的情形，我们选取
\[
 \omega(x,y)=\wp(x-y)dxdy,
\]
其中 $x,y$ 是椭圆曲线的万有覆盖 $\C$ 上的坐标，$\wp(z)$ 是 Weierstrass $\wp$-函数。容易证明 $\operatorname{Res}^2(\omega)=1$，且与 $\omega$ 相伴的射影联络为零。\hfill$\square$

% [Tutor Note: 为什么这两种情形射影联络为零：射影联络 $S_\omega$ 度量的是「$\omega$ 与
%  平坦模型 $dwdz/(w-z)^2$ 的偏离」，而 $\PP^1$ 的 Möbius 变换与椭圆曲线的平移
%  正是使 $\frac{dxdy}{(x-y)^2}$、$\wp(x-y)dxdy$ 保持形式的坐标变换；在这些坐标下
%  上式中的极限恰为零。于是 $T=\tilde T$，联络的曲率只剩 (5.7) 中与 $S_\omega$ 无关的
%  部分，而平移不变性又使该部分为零。]

注意我们的联络确实依赖于双微分 $\omega$ 的选取。下一节我们将更仔细地研究它。这里我们只作为算子来计算其差。设 $\omega'\in H^0(\cC\times_\cB\cC,\omega_{\cC\times_\cB\cC/\cB}(2\Delta))$ 是另一个满足 $\operatorname{Res}^2(\omega)=1$ 的对称双微分。由 (5.5) 与 (5.6) 我们有
\begin{align}
 \nabla^{(\omega)}_X-\nabla^{(\omega')}_X
 &=\big(a_\omega(\vec\ell)-a_{\omega'}(\vec\ell)\big)\cdot\id \nonumber\\
 &=-\frac{c_v}{12}\sum_{j=1}^{N}\operatorname{Res}_{\xi_j=0}
 \big(\ell_j(\xi_j)(S_\omega(\xi_j)-S_{\omega'}(\xi_j))d\xi_j\big)\cdot\id \tag{5.8}
\end{align}
其中 $\theta(\vec\ell)=X$。由于 $S_\omega(z)dz^2-S_{\omega'}(z)dz^2$ 是全纯的，(5.8) 的最后一项与向量场 $X$ 的提升 $\vec\ell$ 的选取无关，并定义一个全纯 1-形式 $\phi_{\omega,\omega'}$，使得
\begin{equation}
 \nabla^{(\omega)}_X-\nabla^{(\omega')}_X=-\frac{c_v}{12}\langle\phi_{\omega,\omega'},X\rangle, \tag{5.9}
\end{equation}
其中 $\langle\ ,\ \rangle$ 是全纯 1-形式与全纯向量场的典范配对。

在本书中，$\cV_{\tilde\lambda}(\mathfrak F)$ 上 $\nabla^{(\omega)}$ 的对偶联络被视为从左作用于 $\cV_{\tilde\lambda}(\mathfrak F)$ 的联络，我们沿用同一记号 $\nabla^{(\omega)}$。于是 $\cV_{\tilde\lambda}(\mathfrak F)$ 上的联络由
\begin{equation}
 \nabla^{(\omega)}_X[|\Phi\rangle]=D(\vec\ell)([|\Phi\rangle])-a_\omega(\vec\ell)[|\Phi\rangle] \tag{5.10}
\end{equation}
给出。

% [Tutor Note: (5.8)–(5.9) 的要点：两个射影平坦联络之差是一个标量 1-形式（乘以 $\id$），
%  即 $\nabla^{(\omega)}-\nabla^{(\omega')}=-\frac{c_v}{12}\langle\phi_{\omega,\omega'},\cdot\rangle$。
%  这个 $\phi_{\omega,\omega'}$ 的存在性正是「射影」二字的分量：若两个联络都平坦，
%  它们的差一般不是一个 1-形式，而是 $\Omega^1\otimes\End$ 中的一个元素；射影平坦
%  保证了差取值于标量。§5.2 将用 $\phi_{\omega,\omega'}$ 把不同 $\omega$ 的扭曲
%  $D$-模粘起来。]

% [Tutor Note: (5.10) 前的这一段原书三处都写作 $\cV_{\tilde\lambda}(\mathfrak F)$（无 $\dagger$），
%  但 $\nabla^{(\omega)}$ 原本是 (5.5) 定义在 $\cV^\dagger_{\tilde\lambda}$ 上的联络，
%  其对偶才作用在 $\cV_{\tilde\lambda}$ 上。按上下文第一处应为
%  $\cV^\dagger_{\tilde\lambda}(\mathfrak F)$；照原文排，见 ISSUES。
%  (5.10) 中 $\cV_{\tilde\lambda}$ 上的联络与 (5.5) 相比多了一个负号，这正是「对偶」的来源：
%  配对 $\langle\Psi|\Phi\rangle$ 求导时两块各贡献一项，符号相反。]

% [原文 p.96]

其中 $\theta(\vec\ell)=X$ 且 $[|\Phi\rangle]\in\cV_{\tilde\lambda}(\mathfrak F)$。这就定义了 $\cV_{\tilde\lambda}(\mathfrak F)$ 上的一个射影平坦联络。$\omega'$ 的不同选取给出如下差：
\begin{equation}
 \nabla^{(\omega)}_X-\nabla^{(\omega')}_X=\frac{c_v}{12}\langle\phi_{\omega,\omega'},X\rangle. \tag{5.11}
\end{equation}
注意这与 (5.9) 的符号相反。

\subsection{扭曲微分算子层}

在上一节中，我们在 $\cV_{\tilde\lambda}(\mathfrak F)$ 与 $\cV^\dagger_{\tilde\lambda}(\mathfrak F)$ 上局部地引入了沿判别轨迹具有正则奇点的可积联络。这些联络只能局部地定义，因为定义联络需要选取满足 (5.2) 的双微分 $\omega$。此外，这些联络依赖于双微分的选取。本节我们将更仔细地研究这种依赖，并定义作用于 $\cV_{\tilde\lambda}(\mathfrak F)$ 与 $\cV^\dagger_{\tilde\lambda}(\mathfrak F)$ 的扭曲微分算子层，它的一阶部分局部地等于联络 $\nabla^{(\omega)}$。

首先我们简要回顾扭曲微分算子层的定义与基本性质。设 $M$ 是复流形，$D=\bigcup_{i=1}^{k}D_i$ 是 $M$ 上的正规交叉除子。

\begin{definition}[5.8]
 若存在 $\cO_\cB$-模与 Lie 代数的正合列
 \[
  0\to\cO_\cB\to\cD_\xi(-\log D)\to\Theta_M(-\log D)\to0,
 \]
 其中 $\cO_\cB$ 被视为由通常乘法诱导的交换 Lie 代数层，则层 $\cD_\xi(-\log D)$ 称为沿 $D$ 具有正则奇点的扭曲微分算子层。
\end{definition}

% [Tutor Note: 定义 5.8 中的 $\cO_\cB$ 应是 $\cO_M$（这里 $M$ 是任意复流形，
%  与底空间 $\cB$ 无关）；原书在此处沿用了上一节的记号，见 ERRATA。]

因此，层 $\cD_\xi(-\log D)$ 作为 Lie 代数层局部同构于沿 $D$ 具有正则奇点的一阶全纯微分算子层 $\cD(-\log D)$。于是，若选取 $M$ 的一个足够细的开覆盖 $\{\cU_\alpha\}_{\alpha\in A}$，我们可以找到平凡化
\[
 \phi_\alpha:\cD_\xi(-\log D)|_{\cU_\alpha}
 \simeq\Theta_{\cU_\alpha}(-\log D)\oplus\cO_{\cU_\alpha},
\]
其中 $\phi_\alpha$ 是 $\cO_\cB$-模同构兼 Lie 代数同态。在 $\cU_{\alpha\beta}=\cU_\alpha\cap\cU_\beta\ne\emptyset$ 上，$\phi_{\alpha\beta}\in\Hom_{\cO_{\cU_{\alpha\beta}}}(\Theta_{\cU_{\alpha\beta}},\cO_{\cU_{\alpha\beta}})$ 由
\[
 \phi_\alpha\circ\phi_\beta^{-1}((X,f))=(X,f+\phi_{\alpha\beta}(X))
\]
定义，其中
\[
 (X,f)\in\Theta_{\cU_{\alpha\beta}}(-\log D)\oplus\cO_{\cU_{\alpha\beta}}.
\]
由下述同构，我们可以把 $\phi_{\alpha\beta}$ 视为 $M$ 上沿 $D\cap\cU_{\alpha\beta}$ 具有对数奇点的亚纯 1-形式：
\[
 H^0\big(\cU_{\alpha\beta},\underline{\Hom}_{\cO_\cB}(\Theta_M(-\log D),\cO_\cB)\big)
 \simeq H^0(\cU_{\alpha\beta},\Omega^1_M(\log D)).
\]
以下我们经常使用这一等同。下面我们证明 $\phi_{\alpha\beta}$ 是 $d$-闭的 1-形式。对 $X,Y\in\Theta_{\cU_{\alpha\beta}}(-\log D)$ 与 $f,g\in\cO_{\cU_{\alpha\beta}}$，由于平凡化是 Lie 代数同态，我们有
\begin{align*}
 [(X,f),(Y,g)]&=([X,Y],X(g)-Y(f))\qquad\text{在 }\cU_\alpha\text{ 上}\\
 &=[(X,f+\phi_{\alpha\beta}),(Y,g+\phi_{\alpha\beta})]\qquad\text{在 }\cU_\beta\text{ 上}\\
 &=([X,Y],(X(g)-Y(f)+X(\phi_{\alpha\beta})-Y(\phi_{\alpha\beta}))).
\end{align*}

% [Tutor Note: 这里的 Lie 代数结构是半直积：$[(X,f),(Y,g)]=([X,Y],X(g)-Y(f))$，
%  其中 $f,g$ 是函数、$X(f)$ 表示向量场作用在函数上。$\cD_\xi(-\log D)$ 的每个
%  截面局部写成 $X+f$（$X\in\Theta(-\log D)$、$f\in\cO$），乘法由这个括号给出。
%  这正是「一阶微分算子 + 零阶项」的代数化：一阶部分 $X$ 与零阶部分 $f$ 的交换子
%  给出 $X(g)-Y(f)$。]

% [Tutor Note: 最后一行的记号有省略：$\phi_{\alpha\beta}$ 在 $X$ 的位置表示
%  $\phi_{\alpha\beta}(X)$、在 $Y$ 的位置表示 $\phi_{\alpha\beta}(Y)$。把两边的第二分量
%  比较（$[X,Y]$ 在 $\cU_\beta$ 下的第二分量还要加上 $\phi_{\alpha\beta}([X,Y])$）即得
%  $X(\phi_{\alpha\beta}(Y))-Y(\phi_{\alpha\beta}(X))-\phi_{\alpha\beta}([X,Y])=0$，
%  也就是 $d\phi_{\alpha\beta}=0$。原书把 $\phi_{\alpha\beta}([X,Y])$ 这一项省掉了。]

% [原文 p.97]

于是有
\[
 \phi_{\alpha\beta}([X,Y])=X(\phi_{\alpha\beta}(Y))-Y(\phi_{\alpha\beta}(X)).
\]
因此我们断言
\[
 d\phi_{\alpha\beta}(X,Y)=X(\phi_{\alpha\beta}(Y))-Y(\phi_{\alpha\beta}(X))
 -\phi_{\alpha\beta}([X,Y])=0.
\]
设 $\Omega^{1,cl}_M(\log D)$ 为沿 $D$ 具有对数奇点的 $d$-闭亚纯 1-形式层；则 $\{\phi_{\alpha\beta}\}$ 是一个余圈，并定义第一上同调群 $H^1(M,\Omega^{1,cl}_M(\log D))$ 中的一个上同调类 $\xi$。

下述命题是以上讨论的容易推论。

\begin{proposition}[5.9]
 沿 $D$ 具有正则奇点的扭曲微分算子层 $\cD_\xi(-\log D)$ 的同构类之集合，与上同调群 $H^1(M,\Omega^{1,cl}_M(\log D))$ 之间存在一一对应。
\end{proposition}

% [Tutor Note: 命题 5.9 的「一一对应」之所以成立，是因为 $H^1$ 恰好分类了
%  局部平凡对象的粘合数据：层 $\cD_\xi(-\log D)$ 由平凡化 $\phi_\alpha$ 与转移
%  1-形式 $\phi_{\alpha\beta}$ 决定，而两个粘合数据给出同构的层当且仅当它们的
%  差是某个 $\{\psi_\alpha\}$ 的上边缘 $\psi_\beta-\psi_\alpha$。这与线丛由
%  $H^1(M,\cO^*_M)$ 分类是同一机制；例 5.10 正是这个双射的具体化：
%  (5.12) 中共轭 $f_\beta^{-1}\circ f_\alpha$ 的零阶部分就是 $\phi_{\alpha\beta}$，
%  对应的类为 $d\log f_{\alpha\beta}$。]

\begin{example}[5.10]
 设 $D=\emptyset$。设 $\cL$ 为 $M$ 上的线丛。存在 $M$ 的开覆盖 $\{\cU_\alpha\}_{\alpha\in A}$，使得有平凡化
 \[
  f_\alpha:\cL|_{\cU_\alpha}\simeq\cO_{\cU_\alpha},
 \]
 并在 $\cU_{\alpha\beta}$ 上置
 \[
  f_{\alpha\beta}:=f_\alpha\circ f_\beta^{-1}.
 \]
 则 $f_{\alpha\beta}$ 是 $\cL$ 的转移函数。与线丛 $\cL$ 相伴的扭曲微分算子层 $\cD_\cL$，是由
 \begin{equation}
  \cD_\cL|_{\cU_\alpha}=f_\alpha^{-1}\circ\cD_{\cU_\alpha}\circ f_\alpha \tag{5.12}
 \end{equation}
 定义在 $\cL$ 上的一阶微分算子层，其中 $\cD_{\cU_\alpha}$ 是 $\cU_\alpha$ 上全纯一阶微分算子层。由 (5.12)，$\cD_\cL|_{\cU_\alpha}$ 到 $\cD_{\cU_\alpha}$ 的 Lie 代数同构由 $f_\alpha$ 给出。因此容易证明，上同调类
 \[
  d\log f_{\alpha\beta}\in H^1(M,\Omega^{1,cl}_M)
 \]
 就是与 $\cD_\cL$ 相伴的那个类。
\end{example}

% [Tutor Note: (5.12) 的共轭写法与上一段的 $\phi_\alpha\circ\phi_\beta^{-1}$
%  互相印证：$f_\alpha$ 把 $\cL$ 上的一阶算子搬成 $\cU_\alpha$ 上的算子，
%  在重叠处两个搬运相差一个共轭 $f_\beta^{-1}\circ f_\alpha$，正是
%  $\phi_{\alpha\beta}$ 的零阶部分。注意 $d\log f_{\alpha\beta}$ 只依赖于转移
%  函数，而 $\cD_\cL\mapsto d\log f_{\alpha\beta}$ 就是命题 5.9 中双射的具体化。]

接下来我们定义 $M$ 上沿 $D$ 具有正则奇点的扭曲微分算子层 $\cD_\xi(-\log D)$ 的特征类。先考虑 $D=\emptyset$ 的情形。存在正合列
\[
 0\to\C\to\cO_M\xrightarrow{\ d\ }d\cO_M\to0.
\]
注意 $d\cO_M=\Omega^{1,cl}_M$。因此由该正合列得到长正合列：
\begin{equation}
 \to H^1(M,\cO_M)\to H^1(M,\Omega^{1,cl}_M)\xrightarrow{\ c\ }H^2(M,\C)
 \to H^2(M,\cO_M)\to. \tag{5.13}
\end{equation}
上述正合列中的映射 $c$ 定义了一个特征类。确切地说，若 $\xi\in H^1(M,\Omega^{1,cl}_M)$ 是对应于 $\cD_\xi$ 的上同调类，则 $c(\xi)\in H^2(M,\C)$ 就是相应的特征类。另有正合列
\[
 0\to\C^*\to\cO^*_M\xrightarrow{\ d\log\ }\Omega^{1,cl}_M\to0.
\]
由该正合列得到长正合列
\begin{equation}
 \to H^1(M,\cO^*_M)\to H^1(M,\Omega^{1,cl}_M)\xrightarrow{\ \bar c\ }H^2(M,\C^*)\to. \tag{5.14}
\end{equation}

% [Tutor Note: 特征类 $c(\xi)$ 的几何含义：$\xi$ 是「粘合 1-形式的类」；
%  若 $\xi$ 可以由 $d\log$ 提升，即来自线丛（如例 5.10），则 $\bar c(\xi)$
%  就是该线丛的第一陈类。$c$ 与 $\bar c$ 分别对应把 $\C$ 系数与 $\C^*$ 系数
%  的指数正合列插入到 de Rham 复形中；$(5.13)$ 与 $(5.14)$ 的差别正是
%  「加性」与「乘性」两种特征类。]

% [原文 p.98]

另一方面，指数映射给出自然的映射
\[
 \exp:H^2(M,\C)\to H^2(M,\C^*),
\]
且我们有等式
\begin{equation}
 \bar c=\exp\circ c. \tag{5.15}
\end{equation}
下述引理是上述正合列 (5.13)、(5.14) 与等式 (5.15) 的容易推论。

\begin{lemma}[5.11]
 设 $D=\emptyset$。$M$ 上的扭曲微分算子层 $\cD_\xi$ 能由例 5.10 借助 $M$ 上的线丛 $\cL$ 给出，当且仅当
 \[
  \bar c(\xi)=0,
 \]
 其中 $\xi\in H^1(M,\Omega^{1,cl}_M)$ 是与层 $\cD_\xi$ 相伴的上同调类。
\end{lemma}

% [Tutor Note: 引理 5.11 说的是：加性特征类 $c(\xi)$ 总能由指数映射变成
%  乘性特征类 $\bar c(\xi)$；$\bar c(\xi)=0$ 意味着 $\xi$ 来自 $\cO^*_M$ 中
%  的余圈（即来自线丛），这正是例 5.10 的情形。所以「扭曲」不是来自
%  线丛时，$\bar c(\xi)\ne0$ 就是障碍。]

当 $D\ne\emptyset$ 时，我们可以类似地定义特征类映射
\[
 c:H^1(M,\Omega^{1,cl}_M(\log D))\to H^2(M\setminus D,\C).
\]
由于后面不会用到这一事实，我们把精确定义留给读者。

现在回到原来的情形。设
\[
 \mathfrak F^{(0)}=(\pi:\cC\to\cB;s_1,\dots,s_N)
\]
是亏格 $g$ 的 $N$ 点稳定曲线的通用族。如 §1.3 所讨论，该族可由通用族 $\mathfrak F^{(0)}=(\mu:\cC^{(0)}\to\cB^{(0)})$ 通过置
\[
 \cB=\underbrace{\cC^{(0)}\times_{\cB^{(0)}}\times_{\cB^{(0)}}\cdots
 \times_{\cB^{(0)}}\cC^{(0)}}_{N}\Big\backslash\bigcup_{i<j}\Delta_{ij},
\]
\[
 \cC=\cC^{(0)}\times_{\cB^{(0)}}\cB
\]
得到，其中 $\Delta_{ij}$ 是第 $(i,j)$ 个对角线，$\pi$ 是到第二个因子的投影。有自然的映射
\[
 \tilde p:\cC\to\cC^{(0)},\qquad p:\cB\to\cB^{(0)}.
\]
注意我们总可以把双微分取作 $\cC^{(0)}\times_{\cB^{(0)}}\cC^{(0)}$ 的某个开集上的双微分的拉回。于是，设 $\{U_\alpha\}_\alpha$ 是 $\cB^{(0)}$ 的开覆盖，使得我们能找到对称双微分
\[
 \omega^{(\emptyset)}_\alpha\in H^0\big(\mu^{-1}(U_\alpha)\times_{U_\alpha}\mu^{-1}(U_\alpha),
 \omega_{\cC^{(0)}\times_{\cB^{(0)}}\cC^{(0)}/\cB^{(0)}}(2\Delta)\big)
\]
满足 $\operatorname{Res}^2\omega^{(\emptyset)}_\alpha\equiv1$。置
\[
 \cU_\alpha=p^{-1}(U_\alpha),\qquad \omega_\alpha=\tilde p^*(\omega^{(\emptyset)}_\alpha).
\]
在 $\cU_\alpha$ 上，我们如下定义 $\Theta_{\cU_\alpha}(-\log D)\oplus\cO_{\cU_\alpha}$ 在 $\cV_{\tilde\lambda}(\mathfrak F)$ 上的作用：
\begin{equation}
 (X,f)\cdot[|\Phi\rangle]=\nabla^{(\omega_\alpha)}_X[|\Phi\rangle]+f\cdot[|\Phi\rangle] \tag{5.16}
\end{equation}
其中 $X\in\Theta_{\cU_\alpha}(-\log D)$、$f\in\cO_{\cU_\alpha}$，而 $[|\Phi\rangle]\in\cV_{\tilde\lambda}(\mathfrak F)$ 是对应于元素 $|\Phi\rangle\in\cH_{\tilde\lambda}(\cB)$ 的元素。（关于 $\nabla^{(\omega_\alpha)}_X$ 的定义见 (5.6)。）在 $\cU_{\alpha\beta}=\cU_\alpha\cap\cU_\beta\ne\emptyset$ 上，两个作用之差由同态
\[
 \varphi_{\alpha\beta}:\Theta_{\cU_\alpha}(-\log D)|_{\cU_{\alpha\beta}}
 \to\cO_{\cU_{\alpha\beta}}
\]
给出。

% [Tutor Note: (5.16) 的关键在于「$f$ 项与 $\omega_\alpha$ 无关」：$\nabla^{(\omega_\alpha)}$
%  本身依赖于 $\omega_\alpha$，但它的一阶部分（$X$ 的方向导数）不变，而零阶部分
%  $a_{\omega_\alpha}$ 已由推论 5.2 证明与 $\omega_\alpha$ 无关。因此 (5.16) 在
%  重叠 $\cU_{\alpha\beta}$ 上只差一个全纯 1-形式 $\varphi_{\alpha\beta}$，它把
%  $\Theta(-\log D)$ 的截面送到函数——这正是定义 5.8 中的半直积结构。]

% [原文 p.99]

由下式定义：
\begin{equation}
 \varphi_{\alpha\beta}(X)=-\frac{c_v}{12}\sum_{j=1}^{N}
 \operatorname{Res}_{\xi_j=0}\big(\ell_j(S_{\omega_\alpha}(\xi_j)
 -S_{\omega_\beta}(\xi_j))d\xi_j\big) \tag{5.17}
\end{equation}
其中 $\theta(\vec\ell)=X\in\Theta_{\cU_{\alpha\beta}}(-\log D)$。现在让我们选取由光滑纤维上的素形式构造的对称双微分 $\omega^{(\emptyset)}_\alpha$。（见 1.1。）于是由 (1.30)，在 $\cU_{\alpha\beta}\setminus D$ 上有
\begin{align*}
 \omega^{(\emptyset)}_\beta(x,y)dxdy&=\omega^{(\emptyset)}_\alpha(x,y)dxdy\\
 &\quad-\frac{1}{2}\sum_{i\le j}\big(v_i(x)v_j(y)+v_i(y)v_j(x)\big)
 \Pi^*\frac{\partial}{\partial\tau_{ij}}\log\det(C\tau+D)dxdy,
\end{align*}
其中 $\Pi$ 是 $U_\alpha$ 到 $\mathfrak S_g$ 的周期映射（见 §1.1.5）。因此我们有
\[
 S_{\omega_\beta}(z)dz^2=S_{\omega_\alpha}(z)dz^2
 -6\Pi^*d\log\det(C\tau+D).
\]
于是在 $U_{\alpha\beta}\setminus D$ 上有
\[
 \varphi_{\alpha\beta}=-\frac{c_v}{2}\Pi^*d\log\det(C\tau+D).
\]
{\emergencystretch=24.2pt 因此，在 $U_{\alpha\beta}\setminus D$ 上 $\varphi_{\alpha\beta}$ 是 $d$-闭的，从而在 $U_{\alpha\beta}$ 上也是 $d$-闭的。于是 $\{\varphi_{\alpha\beta}\}$ 定义了 $H^1(\cB,\Omega^{1,cl}_\cB(\log D))$ 中的一个上同调类。注意，若我们选取其它满足 (5.2) 的双微分 $\{\omega'_\alpha\}$，则得到相同的上同调类 $\{\varphi'_{\alpha\beta}\}=\{\varphi_{\alpha\beta}\}$。也容易证明该上同调类与 $\cB$ 的开覆盖的选取无关。我们把这一上同调类 $[\{\varphi_{\alpha\beta}\}]$ 记为 $\xi(\mathfrak F)$。\par}

% [Tutor Note: (5.17) 的推导链是：$\nabla^{(\omega_\alpha)}-\nabla^{(\omega_\beta)}$
%  由 (5.11) 给出 $= \frac{c_v}{12}\langle\phi_{\omega_\alpha,\omega_\beta},X\rangle$，
%  而 $\phi_{\omega,\omega'}$ 是 $S_\omega-S_{\omega'}$ 的留数配对（见 (5.9) 上方）。
%  于是把 $\ell_j$ 的展开代入即得 (5.17)。第二步用 (1.30)：两个双微分之差是
%  周期映射的拉回再乘上 $\partial_{\tau_{ij}}\log\det$；第三步用
%  $S_\omega$ 的定义（$S_\omega$ 与 $\Pi^*d\log\det$ 的系数关系 $6$）化为
%  $-\frac{c_v}{2}\Pi^*d\log\det(C\tau+D)$。注意 $\varphi_{\alpha\beta}$ 因此
%  正比于 $\Pi^*d\log\det$，故其 $d$-闭性是显然的。]

% [Tutor Note: 定义 5.12 中 $\cD_\cB(-\log D);c_v)$ 的记号在原书里有印刷错误
%  （括号不配对；下标记号也不对）。按上下文应为
%  $\cD_{\xi(\mathfrak F)}(-\log D;c_v)$：下标是上同调类，括号里第二个参数
%  是中心荷 $c_v$。本段之后我们一律写 $\cD_{\xi(\mathfrak F)}(-\log D)$。]

\begin{definition}[5.12]
 扭曲微分算子层 $\cD_\cB(-\log D);c_v)$ 是由命题 5.9 联系于上同调类
 \[
  \xi(\mathfrak F)=\big[\{\varphi_{\alpha\beta}\}\in H^1(\cB,\Omega^{1,cl}_\cB(\log D))\big]
 \]
 的层。
\end{definition}

让我们借助映射类群 $\mathrm{Mod}_g$ 的上同调来描述该上同调类。设 $\tilde\pi:\cC\to\mathbf T_g$ 是 Teichmüller 空间 $\mathbf T_g$ 上亏格 $g$ 光滑曲线的通用族。模群 $\mathrm{Mod}_g$ 作用在该族 $\tilde\pi:\cC\to\mathbf T_g$ 上。遗憾的是，商空间 $M_g:=\mathbf T_g/\mathrm{Mod}_g$ 是奇异的，并且当 $g\ge1$ 时 $M_g$ 上不存在通用族。但存在模群的有限指数子群 $\Gamma_0$，使得通用族 $\tilde{\mathfrak F}^{(0)}:=(\pi:\cC^{(0)}=\tilde\cC/\Gamma_0\to\cB_0=\mathbf T_g/\Gamma_0)$ 存在。注意我们有
\[
 H^i(\mathbf T_g/\Gamma_0,\C)=H^i(\Gamma_0,\C),
\]
\[
 H^i(\mathbf T_g/\Gamma_0,\Omega^{1,cl})=H^i\big(\Gamma_0,
 H^0(\mathbf T_g,\Omega^{1,cl}_{\mathbf T_g})\big),
\]
其中右边把 $\C$ 与 $H^0(\mathbf T_g,\Omega^{1,cl}_{\mathbf T_g}))$ 视为 $\Gamma_0$-模。此外，$M_g=\mathbf T_g/\mathrm{Mod}_g$ 具有 V-流形结构，即只有商奇点的解析空间，并且我们有
\[
 H^i(M_g,\C)=H^i(\mathrm{Mod}_g,\C),
\]
\[
 H^i(M_g,\Omega^{1,cl}_{M_g})=H^i\big(\mathrm{Mod}_g,
 H^0(\mathbf T_g,\Omega^{1,cl}_{\mathbf T_g})\big).
\]
因此我们可以考虑联系于虚拟通用族 $\mathfrak F=(\pi:\cC\to M_g)$ 的扭曲微分算子层的上同调类 $\xi(\mathfrak F)$，或联系于模叠 $M_g$ 上的通用族，即作为 $H^1(\mathrm{Mod}_g,H^0(\mathbf T_g,\Omega^{1,cl}_{\mathbf T_g}))$ 中的元素。

% [原文 p.100]

让我们计算类 $\xi(\mathfrak F)$。对通用族 $\tilde\pi:\tilde\cC\to\mathbf T_g$，利用素形式，我们可以把对称双微分 $\omega(x,y)dxdy$ 定义为层 $p_1^*\omega_{\tilde\cC/\mathbf T_g}\otimes p_2^*\omega_{\tilde\cC/\mathbf T_g}(2\Delta)$ 的截面，使得它在每条纤维上的限制给出满足 $\operatorname{Res}^2\equiv1$ 的对称双微分。（见 (1.38)。）对称双微分的选取依赖于通用族的辛基族的选取。更确切地说，若双微分 $\omega$ 由辛基族 $\{\tilde\alpha_1,\dots,\tilde\alpha_g,\tilde\beta_1,\dots,\tilde\beta_g\}$ 定义，而双微分 $\hat\omega$ 由辛基族 $\{\hat\alpha_1,\dots,\hat\alpha_g,\hat\beta_1,\dots,\hat\beta_g\}$ 定义，则我们有关系
\begin{align}
 \hat\omega(x,y)dxdy&=\omega(x,y)dxdy \tag{5.18}\\
 &\quad-\frac{1}{2}\sum_{i\le j}\big(v_i(x)v_j(y)+v_j(x)v_i(y)\big)
 \frac{\partial}{\partial\tau_{ij}}\log\det(C\tau+D)dxdy, \nonumber
\end{align}
其中
\begin{equation}
 \begin{pmatrix}\hat\beta_1\\ \vdots\\ \hat\beta_g\\ \hat\alpha_1\\ \vdots\\ \hat\alpha_g\end{pmatrix}
 =\begin{pmatrix}A&B\\C&D\end{pmatrix}
 \begin{pmatrix}\tilde\beta_1\\ \vdots\\ \tilde\beta_g\\ \tilde\alpha_1\\ \vdots\\ \tilde\alpha_g\end{pmatrix} \tag{5.19}
\end{equation}
而 $\{\omega_1,\dots,\omega_N\}$ 是该族的标准 1-形式，满足 $\omega_i=v_i(x)dx$。

% [Tutor Note: (5.18) 的实质是：改变辛基（即改变同调基）只改变双微分的
%  周期部分，而差恰好是 $d\tau_{ij}$ 方向的周期映射拉回。$\Pi^*$ 已包含在
%  $\Pi^*(\partial_{\tau_{ij}}\log\det)$ 中；系数 $v_i(x)v_j(y)$ 是标准 1-形式的
%  乘积，故这里把 $\partial_{\tau_{ij}}\log\det(C\tau+D)$ 作为 $\cB$ 上的函数。
%  原书在 (5.18) 第二行的花括号内层用的是圆括号，括号配对系排版错误。]

现在固定辛基族 $\{\tilde\alpha_1,\dots,\tilde\alpha_g,\tilde\beta_1,\dots,\tilde\beta_g\}$ 与该族标准 1-形式的一组基 $\{\omega_1,\dots,\omega_N\}$。于是如 §1.1 那样，我们有周期映射 $\Pi:\mathbf T_g\to\mathfrak S_g$ 与群同态 $\Phi:\mathrm{Mod}_g\to Sp(g,\Z)$，它们具有性质 (1.15) 与 (1.16)。因此，对元素 $\gamma\in\mathrm{Mod}_g$，由 (5.19) 得到新的辛基 $\{\hat\alpha_1,\dots,\hat\alpha_g,\hat\beta_1,\dots,\hat\beta_g\}$，以及新的对称双微分 $\gamma^*(\omega)$。于是由 (5.18) 与 (1.30) 我们有
\begin{align}
 S_{\gamma^*\omega}(z)dz^2&-S_\omega(z)dz^2 \nonumber\\
 &\quad-6\sum_{i\le j}(v_i(z)v_j(z))\Pi^*\Big(\frac{\partial}{\partial\tau_{ij}}
 \log\det(C_\gamma\tau+D_\gamma)\Big)dz^2. \tag{5.20}
\end{align}
于是由命题 1.9 我们有
\[
 S_{\gamma^*\omega}(z)dz^2-S_\omega(z)dz^2
 =6\Pi^*\big(d\log\det(C_\gamma\tau+D_\gamma)\big).
\]
对每个元素 $\gamma\in\mathrm{Mod}_g$ 置
\begin{equation}
 f_\gamma:=\frac{c_v}{2}\Pi^*\big(d\log\det(C_\gamma\tau+D_\gamma)\big). \tag{5.21}
\end{equation}
则 $f_\gamma\in H^0(\mathbf T_g,\Omega^{1,cl}_{\mathbf T_g})$，且容易证明
\[
 f_{\gamma_1\gamma_2}-\gamma_2^*f_{\gamma_1}-f_{\gamma_2}=0.
\]
因此 $\{f_\gamma\}$ 定义了 $H^1(\mathrm{Mod}_g,H^0(\mathbf T_g,\Omega^{1,cl}_{\mathbf T_g}))$ 的一个上闭链。这就是我们的类 $\xi(\mathfrak F)$。$\mathbf T_g$ 上的正合列 (5.14) 被视为 $\mathrm{Mod}_g$-层的正合列，且我们有标准的同态
\[
 H^1(\mathrm{Mod}_g,H^0(\mathbf T_g,\Omega^{1,cl}_{\mathbf T_g}))
 \xrightarrow{\ \tilde c\ }H^2(\mathrm{Mod}_g,\C^*).
\]
% [原文 p.101]

$\xi(\mathfrak F)$ 的像 $\tilde c(\xi(\mathfrak F))\in H^2(\mathrm{Mod}_g,\C^*)$ 由一个二余圈给出：
\begin{align}
 g_{\gamma_1\gamma_2}&=\exp\frac{c_v}{2}\Big\{\log\det(C_{\gamma_3}\Pi(t)+D_{\gamma_3}) \nonumber\\
 &\qquad-\log\det(C_{\gamma_2}\Phi(\gamma_1)\cdot\Pi(t)+D_{\gamma_2})
 -\log\det(C_{\gamma_1}\Pi(t)+D_{\gamma_1})\Big\} \tag{5.22}
\end{align}
其中我们置 $\gamma_3=\gamma_1\cdot\gamma_2$。由于
\[
 \det(C_{\gamma_3}\Pi(t)+D_{\gamma_3})
 =\det(C_{\gamma_2}\Phi(\gamma_1)\cdot\Pi(t)+D_{\gamma_2})
 \times\det(C_{\gamma_1}\Pi(t)+D_{\gamma_1}),
\]
我们断言
\[
 g_{\gamma_1\gamma_2}\in\mu_{2(\ell+g^*)},
\]
其中 $\mu_{2(\ell+g^*)}$ 是 $\C^*$ 的 $2(\ell+g^*)$ 阶循环子群。于是我们得到下述定理。

\begin{theorem}[5.13]
 设我们有位于 Teichmüller 空间 $\mathbf T_g$ 上的通用族之上的亏格 $g$ 光滑曲线的通用族
 \[
  \mathfrak F=(\pi:\cC\to\cB;s_1,\dots,s_N;\eta_1,\dots,\eta_N).
 \]
 则由定义 5.12 联系于扭曲微分算子层 $\cD_\cB(c_v)$ 的上同调类 $\xi(\mathfrak F)$，是由 (5.21) 定义的上同调类
 \[
  \xi\in H^1(\mathrm{Mod}_g,H^0(\mathbf T_g,\Omega^{1,cl}_{\mathbf T_g}))
 \]
 的拉回。特征类 $\tilde c(\xi(\mathfrak F))$ 是由 (5.22) 定义的上同调类
 \[
  \tilde c(\xi)\in H^2(\mathrm{Mod}_g,\mu_{2(\ell+g^*)})
 \]
 的拉回。
\end{theorem}

由于群 $H^2(\mathrm{Mod}_g,\mu_{2(\ell+g^*)})$ 是有限群，而例 5.10 的扭曲微分算子层定义了一个联络，故由定理 5.5 我们得到下述重要的推论。

\begin{corollary}[5.14]
 存在 $\mathrm{Mod}_g$ 的有限指数子群 $\Gamma$，使得若 $\mathfrak F=(\pi:\cC\to\cB;s_1,\dots,s_N;\eta_1,\dots,\eta_N)$ 中的 $\pi:\cC\to\cB$ 由 $\mathbf T_g$ 上的紧 Riemann 曲面族经拉回得到，则共形块层 $\cV^\dagger_{\tilde\lambda}(\mathfrak F)$ 带有射影平坦联络。
\end{corollary}
